\section{Discussion}
\label{sec:discussion}


\subsection{The Latency vs. Flexibility Trade-off}
\label{sec:latency_flexibility_tradeoff}

The empirical measurements establish a clear performance differential. The Seccomp filter denied the \texttt{memfd\_create} system call in 292.82 nanoseconds, while the eBPF LSM program required 832 nanoseconds to deny the subsequent \texttt{fexecve} invocation. This 539.18-nanosecond latency delta represents a direct trade-off between structural restriction and operational flexibility. While this sub-microsecond overhead remains negligible for throughput-oriented web services operating on millisecond-scale latency budgets, it acts as a persistent computational tax when scaled across HPC nodes. Consequently, the analysis must establish the boundary where this latency delta transitions from negligible overhead to a structural performance bottleneck.

Seccomp expends 292.82 nanoseconds on a stateless validation pathway. The kernel evaluates a singular, static condition: whether the incoming system call number matches an entry within a declarative deny list. This binary evaluation maintains no memory of previous invocations and no contextual awareness of the sequence of calls that preceded the current one. Seccomp derives its 292.82-nanosecond execution speed from its stateless architecture.

Because the tracepoint hook attached to \texttt{sys\_enter\_memfd\_create} logs the PID and a nanosecond-resolution timestamp into a BPF hash map, the LSM hook at \texttt{bprm\_check\_security} can subsequently retrieve this entry and evaluate the elapsed temporal delta against a configurable policy window. The program correlates two distinct kernel-level events to map a causal relationship between a memory allocation and an execution attempt. The additional 539.18 nanoseconds yield the capacity to discern behavioral intent.

This distinction introduces direct operational consequences for dynamic, orchestrated cloud environments where, within a Docker Swarm or Kubernetes cluster, container workloads scale continuously while the applied Seccomp profile remains entirely static for the lifecycle of the container instance. Consequently, operators enforcing a global restriction on \texttt{memfd\_create} must accept that any Java Virtual Machine (JVM), Python multiprocessing worker, Node.js rendering process, or build agent utilizing anonymous file descriptors will fail silently during runtime. Static profiles in heterogeneous environments generate high false-positive rates, which ultimately forces administrators to adopt overly permissive security boundaries.

The stateful eBPF implementation mitigates this structural constraint. The correlation logic within \texttt{memfd\_exec\_block.bpf.c} restricts policy enforcement only to the process whose \texttt{comm} field matches the target executable profile, and only when the allocation-to-execution delta falls within the specified detection window. A runtime engineer can tune both parameters independently without modifying application code or restarting container runtimes. The LSM program updates enforcement behavior by reloading the compiled BPF object file via a user-space loader daemon. Kubernetes operators can deliver this update via a DaemonSet rolling update across a thousand-node cluster in the time it would take to rebuild and redeploy a container image with a revised Seccomp JSON profile.

While this 832-nanosecond latency constitutes an acceptable baseline computational cost for architectural adaptability in dynamic cloud environments, it represents a persistent computational tax that degrades cluster throughput when scaled across high-performance computing (HPC) nodes. In environments where compute-bound jobs or Message Passing Interface (MPI) processes perform high-frequency system calls, a 539.18-nanosecond delta accumulates linearly. For example, if a parallel execution thread issues $10^6$ system calls per second, this overhead consumes $0.53918$ seconds of CPU processing time per second, per logical core. This accumulation yields a theoretical 53.918\% throughput degradation. Consequently, this latency penalty stops being negligible and becomes a structural bottleneck. This bottleneck makes stateful eBPF containment incompatible with platforms that demand absolute maximum CPU availability.

\subsection{The Weakness of Structural Defense in HPC Environments}
\label{sec:weakness_structural_defense}

Apptainer's security architecture relies on deterministic workload behavior. High-Performance Computing (HPC) tasks, such as genome sequencing pipelines or finite-element simulations running Message Passing Interface (MPI) jobs, maintain a stable system call surface that does not change between execution cycles. This predictability allows security teams to profile a workflow once and establish a static Seccomp allowlist for the entire cluster. Because scientific software stacks rarely invoke anonymous memory allocation primitives like \texttt{memfd\_create}, the risk of operational false positives remains near zero. This zero-tolerance baseline renders static Seccomp profiles operationally secure.

The container system encodes this Seccomp profile directly at the image level, and the host kernel enforces this structural constraint at the boundary prior to workload execution. The policy travels natively within the Singularity Image Format (SIF) file. This localized design ensures every compute node enforces the identical filter and eliminates dependencies on external monitoring daemons or out-of-band policy synchronization tools.

    However, this architectural approach fails under workload heterogeneity. Introducing a new computational framework, such as a JAX-based machine learning compilation pipeline, shifts this operational profile dynamically. Because operators must first audit the system call surface of the new workflow to update the policy profile, the engineering team is forced to rebuild and redeploy the underlying SIF container image. In ephemeral infrastructure environments where workloads evolve rapidly, managing these static profiles introduces a severe maintenance liability. While the Apptainer Seccomp filter denies the system call in 292.82 nanoseconds without active daemon overhead, the 832-nanosecond eBPF framework provides the runtime adaptability necessary to accommodate heterogeneous workloads without continuous profile regeneration. 

\clearpage 

\subsection{Architectural Recommendations}
\label{sec:architectural_recommendations}

The divergence between dynamic cloud orchestration and deterministic high-performance computing requires a workload-specific security architecture. To address the Scanner Gap across standard enterprise infrastructures, security architectures deploy an environment-specific defensive taxonomy. This taxonomy divides defenses into two operational categories based on runtime volatility.

For deterministic, static infrastructures, security frameworks implement structural containment at the image boundary. Apptainer environments running coupled scientific workloads enforce strict containment by applying default-deny Seccomp profiles during container compilation. Because the workload profile remains immutable during its lifecycle, stripping capabilities and blacklisting system calls such as \texttt{memfd\_create} introduces zero operational overhead. This static configuration eliminates the performance tax of active host daemons. By forcing the security barrier down into the kernel entry path, this approach preserves throughput and prevents the exposure of execution paths to runtime manipulation.

Heterogeneous enterprise platforms, such as multi-tenant Kubernetes clusters, experience continuous microservice churn and unpredictable third-party application behaviors. To secure these environments, operators deploy eBPF-based Linux Security Modules (LSM) via cluster-wide agents to establish stateful observability layers. By evaluating stateful system call correlation windows, the infrastructure retains the capacity to allow legitimate execution primitives. The underlying monitoring agent terminates correlated memory injection vectors at the LSM boundary. This dynamic deployment model provides the engineering flexibility necessary to secure multi-tenant microservice environments. Ultimately, the empirical data proves that a unified defensive posture is impossible; security engineers must bifurcate their runtime defenses based strictly on the system call volatility of the workload, selecting either low-latency structural containment for static tasks or stateful, adaptive correlation for dynamic services.